\begin{slide}
\begin{center}\Title{Our proposal}\end{center}

\begin{itemize}

\IncPart{1}{\item[]
We recall that $\Uniform{h}\Appl{n} \bydef n + h$
appears first in early de Bruijn's works,
where the theory of update functions finds its foundation.
}

\IncPart{2}{\item[]
A class of systems originating from these works and from $\LS$
\stress{delay} the application of update functions
by storing them in their syntax.
}

\IncPart{3}{\item[]
The main reason for so doing is allowing
a computation either to delay updates at will or to apply them whenever is the case.
}

\IncPart{4}{\item[]
The system we are presenting here today belongs to this class.
}

\IncPart{5}{\item[]
We wish to introduce a $\beta$-reduction step $\EFStep{}$
that is \stress{focused} and \stress{expanding},
We explain these terms by means of an example.
}

\IncPart{5}{\item[]
\begin{center}
$\Comp{\Abst{x}x\Appl{x}}\Appl{V} =
 \Comp{\Abst{}1\Appl{\mstress{1}}}\Appl{V} \EFStep{}
 \Comp{\Abst{}1\Appl{\mstress{\IRef{1}V}}}\Appl{V}$
\end{center}
}

\IncPart{5}{\item[]
The step attaches a copy of $V$ after the leaf $\mstress{1}$,
which becomes an inner node
and records that we still need to update $V$ with
$\Uniform{\mstress{\scriptstyle 1}}$.
}

\IncPart{6}{\item[]
\stress{Focused} means we replace exactly one instance of $x$.\\
\stress{Expanding} means the redex is a subtree of the reduct.
}

\end{itemize}

\end{slide}
